Praktika honen helburua robot mugikor baten sentsoreekin lan eginez bi portaera autonomo garatzea da. Lehenengoan, robotak ingurune batean nabigatu behar du oztoporik jo gabe, LiDARraren informazioa bakarrik erabiliz. Bigarrenean, errepide-zirkuitu bati jarraitu behar dio kanpoaldetik, aurreko kameraren irudian oinarrituta.

Garatutako algoritmoak ROS~2 erabiliz inplementatu dira Gazebo simulagailuan, eta lan osoa ROSMASTER X3 robotarentzat diseinatu da, bere bi sentsore nagusiak erabiliz: RPLIDAR S2 LiDARra eta aurrealdeko RGB-D kamera. Errepide-jarraipena, gainera, robot errealean ere ebaluatu dugu.
